% Original reconstruction of source Fig.90.1, PDF372 / printed360.
% This is a basic-TikZ input fragment; the standalone PDF is generated from it.
% P is the physical incoming pion momentum; k1,k2 are physical outgoing photons.
% The axial current kernel follows S76 with all momenta outgoing: r=-P.
% Black arrows on solid lines are the fermion/momentum directions of the
% selected loop routing. Grey photon arrows indicate outgoing physical momenta.
% First orientation:
%   upper -> axial: ell-k1; lower -> upper: ell; axial -> lower: ell+k2.
% Reversed orientation:
%   axial -> upper: ell+k1; upper -> lower: ell; lower -> axial: ell-k2.
% These correspond to the two traces described in MW.07--12 of
% checks/review_S90_massive_ward_independent.md.
% All three dots are local current insertions. No vertex factor is assigned.
% The unexplained infinity-shaped marks in the source scan are not reproduced
% and are not assigned a mathematical meaning.
\begin{tikzpicture}[x=1mm,y=1mm,line width=.7pt,>=latex,
  line cap=round,line join=round,
  font=\fontsize{10.5}{12}\selectfont,
  momentum/.style={->,line width=.45pt,draw=black!58}]
  \foreach \shift/\orientation in {0/1,80/2} {
    \begin{scope}[xshift=\shift mm]
      % Fixed original layout shared by the two opposite loop orientations.
      \draw (0,0)--(27,16)--(27,-16)--cycle;
      \draw[dashed] (-18,0)--(0,0);
      \draw[->] (-13,0)--(-6,0);
      \node[anchor=west,inner sep=0pt] at (-18,4) {$\pi^0$};
      \node[anchor=south,inner sep=0pt] at (-8,2) {$P$};

      % Two sinusoidal vector photon paths; separate arrows fix k1,k2 outward.
      \draw plot[domain=0:1,samples=141,variable=\t]
        ({27+18*\t-.30*sin(360*5*\t)},
         {16+8*\t+.68*sin(360*5*\t)});
      \draw plot[domain=0:1,samples=141,variable=\t]
        ({27+18*\t+.30*sin(360*5*\t)},
         {-16-8*\t+.68*sin(360*5*\t)});
      \draw[momentum] (34,22.5)--(43,26.5);
      \draw[momentum] (34,-22.5)--(43,-26.5);
      \node[anchor=south,inner sep=0pt,text=black!65] at (38.5,28) {$k_1$};
      \node[anchor=north,inner sep=0pt,text=black!65] at (38.5,-28) {$k_2$};

      \fill (0,0) circle[radius=.5mm];
      \fill (27,16) circle[radius=.5mm];
      \fill (27,-16) circle[radius=.5mm];
      \node[anchor=west,inner sep=0pt] at (29.5,0) {$\ell$};

      \ifnum\orientation=1
        \draw[->] (18,10.666667)--(9,5.333333);
        \draw[->] (27,-5)--(27,5);
        \draw[->] (9,-5.333333)--(18,-10.666667);
        \node[anchor=south,inner sep=0pt] at (11,11.5) {$\ell-k_1$};
        \node[anchor=north,inner sep=0pt] at (11,-11.5) {$\ell+k_2$};
        \node[inner sep=0pt] at (11,27) {$(a)$};
      \else
        \draw[->] (9,5.333333)--(18,10.666667);
        \draw[->] (27,5)--(27,-5);
        \draw[->] (18,-10.666667)--(9,-5.333333);
        \node[anchor=south,inner sep=0pt] at (11,11.5) {$\ell+k_1$};
        \node[anchor=north,inner sep=0pt] at (11,-11.5) {$\ell-k_2$};
        \node[inner sep=0pt] at (11,27) {$(b)$};
      \fi
    \end{scope}
  }
  \node[inner sep=0pt] at (54,0) {$+$};
  \node[inner sep=0pt] at (53.5,-37)
    {$P=k_1+k_2,\qquad r=-P\quad\text{（轴核全出）}$};
\end{tikzpicture}

